\section{Law charts and the height-two quotient}
\label{geom:section}

The counterexample will be built in a category of periodic complexes
of sheaves on a stack of formal group laws.  Our task here is to
construct the stack and its chart sheaf, then establish tools for
computing sections.  We begin with the integral stack of laws of height
at most two, pass to characteristic $p$, and quotient a central
procyclic group of automorphisms.  The quotient has a simple height-one
open and explicit nilpotent neighborhoods of height two.  The height-two
stabilizer provides both the coefficient ring for the construction and
the finite cohomological bounds needed for periodic complexes.

Throughout this section $p>5$ is prime.  Put
\[
 D=2(p-1),\qquad m_i=\frac{p^i-1}{p-1},\qquad q=m_2=p+1.
\]
Additional large-prime conditions will be imposed only when constructing
the operator used in the counterexample.

\subsection{The integral law groupoid}

We start over $\Z_{(p)}$ because the same flat chart will supply the
characteristic-$p$ construction and the later integral comparison.
We use the classical Brown--Peterson Hopf algebroid
\[
 P=\Z_{(p)}[v_i:i\geq1],\qquad
 \Gamma_P=P[t_i:i\geq1]
\]
of strict $p$-typical formal group laws.  Divide all its topological
degrees by $D$, so that $v_i$ and $t_i$ have weight $m_i$.
The formal-group and cooperation calculations used here are recalled
in \cite[Chapter~4]{Ravenel1986}; for the Araki formulas see also
\cite[Appendix~B.5.5--B.5.6, B.5.12, and B.5.15]{Ravenel1992}.
We choose Araki coordinates, for which
the universal law $F$ satisfies
\begin{equation}\label{geom:p-series}
 [p]_F(T)=\mathop{\textstyle\sum}\nolimits_{i\geq0}^{F}
             v_iT^{p^i},\qquad v_0=p.
\end{equation}
The superscript denotes addition in $F$, not ordinary addition.
The following calculation recalls why this is an integral coordinate
choice.  Write $\log_F(T)=\sum_{i\geq0}l_iT^{p^i}$ and successively
subtract the lower terms of \eqref{geom:p-series} from $[p]_F$ using
formal-group subtraction.  The logarithm shows that the next remainder starts at
$T^{p^n}$; its coefficient is integral, since all the preceding
subtractions were integral.  Comparison with the Hazewinkel recursion
\[
 p l_n=v_n^{\mathrm{Haz}}
       +\sum_{0<i<n}l_i(v_{n-i}^{\mathrm{Haz}})^{p^i}
\]
gives a triangular change of generators with diagonal
$1-p^{p^n-1}$, a unit in $\Z_{(p)}$.

A strict arrow from $F_R$ to $F_L$ has the form
\[
 \phi(T)=T+_{F_L}\mathop{\textstyle\sum}\nolimits_{i\geq1}^{F_L}
                      t_iT^{p^i}.
\]
The coefficients of $F_R$ are given by the right unit $\eta_R$.
Commuting $\phi$ with the $p$-series gives
\begin{equation}\label{geom:invariant-ideals}
 \eta_Rv_i\equiv v_i\pmod{(p,v_1,\ldots,v_{i-1})}.
\end{equation}
In particular the ideal $(p,v_1,v_2)$ defines an invariant closed locus.

Let $\mathcal M$ be the stack of this groupoid on
$\Spec P\setminus V(p,v_1,v_2)$, also divided by $\Gm$ for the
regraded weights.  All stacks and sheaves in this section use fpqc
descent.  Define
\[
 \begin{aligned}
 B_2&=\Z_{(p)}[v_1,v_2],\\
 A_2&=B_2[v_2^{-1}],\\
 \mathcal X&=[(\Spec B_2\setminus V(p,v_1,v_2))/\Gm],
 \end{aligned}
 \qquad f:\mathcal X\longrightarrow\mathcal M,
\]
where $f$ puts $v_{i>2}=0$.  We will also use the standard
Johnson--Wilson coefficient chart
\[
 A_{\mathrm{Haz}}
 =\Z_{(p)}[v_1^{\mathrm{Haz}},(v_2^{\mathrm{Haz}})^{\pm1}],
\]
with the higher Hazewinkel generators set equal to zero.

\begin{proposition}\label{geom:law-charts}
The map $f$ is flat and affine.  The maps
$\Spec A_2\to\mathcal M$ and
$\Spec A_{\mathrm{Haz}}\to\mathcal M$ are flat affine covers.
In graded coordinates, the cooperation Hopf algebroid of the latter
chart has arrow ring
\[
 A_{\mathrm{Haz}}\otimes_{P,\eta_L}\Gamma_P
                  \otimes_{P,\eta_R}A_{\mathrm{Haz}}.
\]
\end{proposition}

\begin{proof}
We verify flatness, affineness, and then coverage.  The Landweber
flatness criterion applies to each localization of $B_2$
inverting $p$, $v_1$, or $v_2$, and to $A_{\mathrm{Haz}}$: the invariant
regular sequence remains regular up to the point where its ideal
becomes the unit ideal.  The interpretation of this criterion as
flatness of law charts is the standard equivalence between Landweber
exact comodules and flat descent; see
\cite[Lemma~2.2, Corollary~2.3, and Theorem~2.5]{HS2005}
and \cite[Section~5, Theorem~26, and Proposition~27]{Naumann2007}.

One can pass explicitly from exactness on graded comodules to the
flatness assertion being used.  Apply exact coinduction
$\Gamma_P\otimes_P(-)$ to a sequence of graded $P$-modules, and then
apply base change to the chart ring $A$.  This tests flatness of
$A\otimes_{P,\eta_L}\Gamma_P$ over $P$ via $\eta_R$; inversion of
arrows interchanges the two units if necessary.  Graded flatness here
implies ordinary flatness.  For an ungraded test module $N$, put a copy
of $N$ in every weight and let each homogeneous element act through its
ordinary action into the correspondingly shifted copy.  A fixed
weight of the graded tensor product recovers the ungraded tensor
product.  Thus the graded exactness test also tests arbitrary modules.

After pulling back to an affine scheme of laws with its grading frame
trivialized, the fiber of $f$ is the affine scheme of strict arrows
whose other end is a $B_2$-law.  The deleted locus cannot occur at that
end by \eqref{geom:invariant-ideals}.  This proves affineness of $f$.
Fiber products of affine law charts are affine as well: in addition
to the strict-arrow coordinates, they retain only an invertible
grading parameter.  In particular the relevant chart maps and
diagonals are affine.

It remains to check that the indicated flat charts cover.  This may
be checked on geometric points.  In characteristic zero the logarithm
gives strict isomorphisms of the required $p$-typical form.  In
characteristic $p$, let $h=1$ or $2$ be the height, and compare two
laws over an algebraically closed field with the same nonzero
coefficient $v_h$.  For a fixed left law, solve for an arrow by
prescribing the right coefficients in order.  After the preceding
equations have been imposed, the equation at step $j\geq1$ is
\begin{equation}\label{geom:recursive-arrow}
 \eta_Rv_{j+h}
 =v_h t_j^{p^h}-v_h^{p^j}t_j
       +\text{a polynomial already determined by }t_{<j}.
\end{equation}
Indeed, compare degree $p^{j+h}$ in
$[p]_{F_L}(\phi(T))=\phi([p]_{F_R}(T))$.  All exponents in the
left $p$-series are divisible by $p^h$, so only $\phi$ through degree
$p^j$ can enter there.  The first occurrence of $t_j$ is additive in
degree $p^j$.  On the right, apart from the displayed coefficient
$\eta_Rv_{j+h}$, only previously fixed coefficients enter: the
formal-sum coefficients needed at this step have total exponent at
most $p^j$ and hence, by homogeneity, use only generators of weight
at most $m_j$.  This gives \eqref{geom:recursive-arrow}.
Its leading coefficient is nonzero, so it can be solved at every
step.  Both charts contain laws of the prescribed height and leading
coefficient; for the Hazewinkel chart use the homogeneous triangular
coordinate change above.  This proves coverage.  Finally the displayed
arrow ring is the usual Landweber base-change calculation for
cooperations, and is also the ring obtained directly by pulling back
the two ends of the law groupoid.
\end{proof}

We fix the chart sheaves
\begin{equation}\label{geom:integral-chart-sheaves}
 T_{\Z}=f_*\mathcal O_{\mathcal X},\qquad
 T_{\Z_p}=f_*\bigl(\mathcal O_{\mathcal X}
                          \otimes_{\Z_{(p)}}\Z_p\bigr).
\end{equation}
They are $p$-monic, meaning that multiplication by $p$ is injective.
This follows on $\mathcal X$ and is preserved by the affine pushforward.
The later comparison will identify these particular sheaves with the
algebraic images of Brown--Peterson spectra.

\subsection{Removing the central procyclic group}

The construction itself uses the characteristic-$p$ law stack.  On this
stack a central copy of $\Z_p$ acts through strict automorphisms of
every law.  We remove this common inertia to obtain a height-two
stabilizer whose invariants have cohomological dimension at most three.

Put $k=\F_{p^2}$.  Write
\[
 i:\mathcal M_0\longrightarrow\mathcal M,
 \qquad e:\mathcal M_k\longrightarrow\mathcal M_0
\]
for reduction to characteristic $p$ and extension of scalars from
$\F_p$ to $k$, respectively.  Subscripts $0$ and $k$ on $\mathcal X$
and $f$ mean the corresponding base changes.
The proconstant group
\[
 Z=1+p\Z_p\simeq\Z_p
\]
acts as central inertia by the strict automorphisms $[z]_F$.
Here a proconstant group scheme is an inverse limit of finite
constant group schemes; its functions over $k$ are locally constant
$k$-valued functions.  Integer multiplication has the $p$-typical
formal-sum expansion obtained from logarithms and successive
subtraction as above.  In characteristic $p$, the $p$-series has
formal-variable order at least $p$, so these multiplications extend
continuously to $p$-adic integers on every jet.  They commute with
all arrows and with the grading.  The action is faithful on geometric
fibers because the laws have finite height.  Fix a topological
generator $c$ of $Z$.

We construct an affine groupoid presentation for the quotient and
prove that the chart remains affine.  These properties will make its
pushforward accessible to the later coinduction calculation.

\begin{lemma}\label{geom:torsor}
Let $A$ be a $k$-algebra with a locally finite continuous action of
$Z$, and let $Z_m$ be the subgroup of index $p^m$.  Suppose that
\[
 A=\bigcup_{m\geq0}A^{Z_m}
\]
and that the action on $\Spec A$ is geometrically free.  Then
$\Spec A\to\Spec(A^Z)$ is an fpqc $Z$-torsor.
\end{lemma}

\begin{proof}
We first descend the free action to finite quotients.  We then prove
the finite torsor assertion directly and pass to the limit.
Set $A_m=A^{Z_m}$.  Every element has finite orbit, so $A$ is
integral over each $A_m$.  Two geometric points with the same
restriction to $A_m$ differ by $Z_m$.  To see this, restrict first
to a subalgebra generated by finitely many finite orbits.  For a finite
group, distinct orbits of geometric points are separated by invariant
functions: apply the Chinese remainder theorem and take an invariant
product of separating functions.  Formation of finite-group
invariants commutes with field extension, since it is a finite
kernel and a field extension is flat.  The finite-level assertions
then give nonempty closed sets of possible elements of $Z_m$;
compactness supplies a compatible element for the whole algebra.

Integrality lifts every geometric point of $A_m$ to one of $A$.
The induced $Z/Z_m$-action on $\Spec A_m$ is free.  In fact, if a
coset fixes such a point, lift it to a point of $A$.  The translate
of that lift and the original lift differ by $Z_m$ by the preceding
paragraph.  Geometric freeness upstairs forces the original coset
to be trivial.

For completeness, a finite constant group $K$ acting geometrically
freely on an affine scheme gives a torsor over the invariant ring
even when $|K|$ is not invertible.  At a lift of a geometric point
of the invariant spectrum, choose elements $b_1,\ldots,b_{|K|}$
whose evaluations on its orbit form an invertible matrix.  Such
elements exist by the Chinese remainder theorem after field extension,
and hence a determinant of elements from the original algebra is
nonzero there.  If
$D_b=\det(g(b_j))_{g\in K,j}$, then $D_b^2$ is invariant.
After inverting it, solve for the coefficients of any element in
the basis $b_j$ using this matrix.  Uniqueness makes the coefficients
invariant.  The same matrix gives the tensor torsor isomorphism.
These neighborhoods cover the invariant spectrum by integrality.

Applying this argument to $A_m$ shows that $\Spec A_m$ is a finite
faithfully flat $Z/Z_m$-torsor over $\Spec A^Z$.  Passing to the
filtered union gives faithful flatness and the torsor identity for
$A$.  The resulting affine map is quasi-compact, hence fpqc.
\end{proof}

\begin{proposition}\label{geom:rigidification}
There is a flat affine groupoid stack $\mathcal N$ and a map
$r:\mathcal M_k\to\mathcal N$ that quotients the vertical central
$Z$-action on arrows.  Its pullback is exact and identifies sheaves
on $\mathcal N$ with sheaves on $\mathcal M_k$ having trivial vertical
$Z$-action.  Moreover
\[
 f'=rf_k:\mathcal X_k\longrightarrow\mathcal N
\]
is flat and affine.  The weight line descends to $\mathcal N$, and
$t=v_1$ is an invariant Cartier section of $\mathcal O(1)$ there.
\end{proposition}

\begin{proof}
We construct the quotient on an affine presentation, and then check the
chart and the weight line.  Take $S=\Spec(A_2\otimes k)$ and quotient
its affine arrow scheme
$S\times_{\mathcal M_k}S$ by vertical composition with $Z$.
The hypotheses of Lemma~\ref{geom:torsor} hold.  The action is
geometrically free by the leading orders of finite-height $p$-series,
and it is locally finite on the arrow coefficients by uniform
continuity on each formal-variable jet.  Thus the quotient is affine
and the original arrow scheme is an fpqc torsor over it.
Composition descends by centrality: after locally lifting two
composable arrows, changing either lift changes their composite by
the corresponding vertical element.  Identity and inverse descend
as well.  Source and target remain flat by faithfully flat descent.
This defines the asserted groupoid and $\mathcal N$.
Descent of a sheaf's arrow isomorphism through the torsor gives the
claimed description of $r^*$ and its exactness.

Similarly, $S\times_{\mathcal N}\mathcal X_k$ is the quotient of
$S\times_{\mathcal M_k}\mathcal X_k$ by vertical composition.
An isomorphism in the quotient has local lifts whose choices form
a $Z$-torsor, as is visible on the arrow chart.  There is no residual
automorphism from $\mathcal X_k$ in this vertical subgroup: an
automorphism from $\mathcal X_k$ is a grading element, and its
grading parameter must be identity to lie in the strict subgroup.
This assertion holds fpqc locally over arbitrary bases, not only on
geometric points.  The quotient is therefore represented by the
invariant affine of Lemma~\ref{geom:torsor}.  Its map to $S$ is flat,
since the map upstairs is flat and the torsor is faithfully flat.
This proves the assertion about $f'$.

The torsor descriptions commute with base change.  The central
action is trivial on the weight line, so that line descends.
Equation~\eqref{geom:invariant-ideals} in characteristic $p$ makes
$v_1$ invariant.  It is a non-zero-divisor on the law charts, proving
the Cartier assertion.
\end{proof}

Use the same notation $t$ on $\mathcal M_0$ and $\mathcal M_k$, and set
\begin{equation}\label{geom:characteristic-chart-sheaves}
 T_0=f_{0*}\mathcal O_{\mathcal X_0},\qquad
 T_k=f_{k*}\mathcal O_{\mathcal X_k},\qquad
 T'=f'_*\mathcal O_{\mathcal X_k}.
\end{equation}
The settings and their chart objects are arranged as follows:
\[
\begin{array}{ccccccc}
 \mathcal N&\xleftarrow{r}&\mathcal M_k&\xrightarrow{e}&\mathcal M_0&
 \xrightarrow{i}&\mathcal M\\
 T'&&T_k&&T_0&&T_{\Z},\,T_{\Z_p}.
\end{array}
\]
The upper row records maps of stacks: $r$ quotients the central
automorphisms, $e$ extends scalars from $\F_p$ to $k$, and $i$ is the
characteristic-$p$ inclusion.  The lower row records the chart objects
in those settings.  Sections~\ref{complete:section}--\ref{obs:section}
carry out the algebraic construction on $\mathcal N$ in the periodic
category defined next, with $T'$ as its chart object.
Sections~\ref{comp:section}--\ref{transfer:section} return to the
integral setting and transfer the construction to spectra.

The Cartier section $t$ determines the height-one open and the
height-two thickenings.  We next specify the category and the section
groups in which the construction is measured.

\subsection{Periodic complexes and sections}
\label{geom:periodic-foundations}

Later we will assemble quotients by powers of $t$ using inverse limits,
so the category must allow arbitrary quasi-coherent sheaves as well as
finite constructions.  We fix its conventions before describing the
charts on which those quotients are computed.  The weight twist of a
graded module is
\[
 W(s)_a=W_{s+a}.
\]
For a stack $\mathcal S$ with a flat affine groupoid presentation and
weight line $\mathcal O(1)$, let $\mathbf D(\mathcal S)$ be the derived
category of twisted periodic homological complexes of quasi-coherent
sheaves:
\begin{equation}\label{geom:periodic-complex}
 C_{j+D}=C_j(1),\qquad d:C_j\longrightarrow C_{j-1}.
\end{equation}
The differential respects the periodic identifications, and we invert
quasi-isomorphisms.  We retain the associated stable homotopy theory
when forming homotopy limits or colimits.  Write $C[h]=\Sigma^hC$,
with $(C[h])_j=C_{j-h}$ and the usual sign on the differential.
A sheaf $W$ denotes the pure complex with $C_0=W$ and zero terms
outside the residue class $0\pmod D$, extended by the twists.

For the integral stack $\mathcal M$, Section~\ref{comp:section}
identifies this periodic category, by flat descent and an explicit
degree regrading, with the differential $E_*E$-comodule model used in
the chromatic comparison, where $E=E(2)$.

For any of these categories set
\begin{equation}\label{geom:sections-homotopy}
 \pi_j C=[\mathcal O[j],C]_{\mathbf D(\mathcal S)}.
\end{equation}
This is sections homotopy, not underlying sheaf homology.  For example,
a sheaf map $W\to W'(s)$ extends to a map $W[sD]\to W'$ of pure
complexes.  Such sheaf maps need not account for all derived maps.
After the chart calculations, we will construct the model structure
and finite resolutions that justify our computations of these derived
groups.

\subsection{The height-one open and height-two thickenings}

We now describe the two pieces determined by $t$.  The open has no
remaining stabilizer after the central quotient.  On each nilpotent
thickening, the stabilizer acts on one explicit coordinate ring; this
will turn the reductions of $T'$ into concrete coefficient modules.

\begin{proposition}\label{geom:open}
There are equivalences
\[
 \mathcal M_0[1/t]\simeq BZ,\qquad
 \mathcal N[1/t]\simeq\Spec k.
\]
The first is split by the law with $v_1=1$ and $v_{i>1}=0$.
\end{proposition}

\begin{proof}
Normalize $v_1=1$ by its weight-one grading.  We will show that every
normalized law is locally isomorphic to the stated law and then
determine its stabilizer.  An arrow from the stated law into an arbitrary
normalized law is obtained by imposing
$\eta_Rv_{j+1}=0$ successively.  Modulo the preceding equations,
\[
 \eta_Rv_{j+1}=t_j^p-t_j
                    +\text{a polynomial in }t_{<j}.
\]
This is the height-one case of the coefficient comparison in
\eqref{geom:recursive-arrow}; equivalently, only $\phi^p$ on the
left and $\phi(T^p)$ on the right can contribute the new $t_j$.
The equations are successively finite \'etale, so their limit is
faithfully flat and supplies the cover.

On the splitting law the automorphism equations reduce recursively
to $t_j^p=t_j$.  Their group of solutions is exactly $Z$:
$[1+ap^j](T)-_F T$ supplies any prescribed next digit in $\F_p$,
and successive correction and completeness identify all strict
automorphisms.  There is no remaining grading stabilizer, since
strict arrows preserve the normalized coefficient $v_1=1$.
The first equivalence follows.  Base change to $k$ and quotient by
the central group give the second.
\end{proof}

Let
\[
 \Lambda=\mu_q(k),\qquad R_d=k[x]/(x^d)\quad(d\geq1),
 \qquad R=k[[x]].
\]
The group $\Lambda$ has order $q$, prime to $p$.  The rings $R_d$
will be the coordinate rings of the thickenings, and $R$ records their
compatible coefficients.  Let $F_0$ be the Honda law with $v_2=1$
and $v_i=0$ for $i\ne2$, $i\geq1$.
Denote its proconstant group of strict automorphisms by $G^1$, and put
\[
 G_{\mathrm{up}}=G^1\rtimes\Lambda,\qquad
 H=G^1/Z,\qquad G=H\rtimes\Lambda.
\]
The action of $\Lambda$ on strict automorphisms is the weight action.

\begin{proposition}\label{geom:thickened-charts}
The chart law $F_x$ with $v_1=x$, $v_2=1$, and $v_{i>2}=0$
gives compatible presentations
\begin{equation}\label{geom:quotient-thickenings}
 \mathcal M_k/(t^d)=[\Spec R_d/G_{\mathrm{up}}],\qquad
 \mathcal N/(t^d)=[\Spec R_d/G].
\end{equation}
Under these presentations, $\mathcal X_k/(t^d)$ is
$[\Spec R_d/\Lambda]$.  Sheaves are smooth, that is, discrete
continuous, semilinear representations over $R_d$.
The compatible coefficient action on $R$ satisfies
\[
 g(x)\in x+x^2R\quad(g\in H),\qquad
 \lambda(x)=\lambda^{-1}x\quad(\lambda\in\Lambda).
\]
If $u$ is the chart basis of the weight line with $t=xu$, then
\begin{equation}\label{geom:twist-cocycle}
 g(u)=\frac{x}{g(x)}u\quad(g\in H),\qquad
 \lambda(u)=\lambda u\quad(\lambda\in\Lambda).
\end{equation}
For every strict label representing an element $g\in H$, there is
a formal isomorphism $F_{g(x)}\to F_x$ over $R$ reducing to that label.
\end{proposition}

\begin{proof}
We normalize $v_2$ and then identify arrows by lifting their
special-fiber labels.  Uniqueness of the lifts will also identify
composition over every nilpotent base.

On $V(t^d)$ the coordinate $v_2$ is invertible.  Normalize it to 1
by adjoining a $q$-th root of the grading parameter.  This gives a
flat cover by the stated chart: the map obtained by adjoining a
grading parameter $\lambda$ has coordinates
$(v_1,v_2)=(\lambda x,\lambda^q)$ and is the composite of the
original flat cover with the grading action.  On $\mathcal X_k$ the
remaining grading group is $\Lambda$, giving the asserted chart
there; its ambient $v_2\ne0$ open is $[\Spec k[x]/\Lambda]$.

We identify all arrows and their composition, not just their
geometric fibers.  If $\rho$ is the grading parameter and $t_j$ are
the strict-arrow parameters with left end $x$, the right chart point
is $y=x/\rho$, since $\eta_Rv_1=x$.  The remaining equations are
\begin{equation}\label{geom:nilpotent-arrow-equations}
 \eta_Rv_2=\rho^q,\qquad \eta_Rv_{j+2}=0\quad(j\geq1).
\end{equation}
Modulo $x$ the first right unit is 1.  Modulo the preceding equations
the subsequent ones are $t_j^{p^2}-t_j=0$.  This follows by comparing
$\phi(T)^{p^2}$ with $\phi(T^{p^2})$ for the Honda law; its
coefficients lie in $\F_p$, and all preceding digits are fixed by
the square of Frobenius.  Thus the special-fiber arrow scheme is
proconstant with labels
\[
 (\rho_0,t_1^0,t_2^0,\ldots)
       \in\Lambda\times k^{\mathbb N_{>0}}.
\]
Labels with $\rho_0=1$ give $G^1$, and grading-only arrows split
off $\Lambda$.

For a prescribed left end $x$, each label lifts uniquely through
the nilpotent layers.  At a square-zero layer
$(x^b)/(x^{b+1})$, only the special-fiber Jacobian acts.  Its first
diagonal entry is a unit because $q$ is prime to $p$; the other
diagonal entries are $-1$.  After adding combinations of earlier
rows, the system is lower triangular.  Each equation is a finite
polynomial, so the linear correction system on the infinite product
of coefficient modules is solved recursively from its first row.
This proves existence and uniqueness on every layer and over every
chart algebra.  Proconstant labels are unchanged by a nilpotent
thickening, so the construction is functorial.  In particular
composition is determined by composition of the special-fiber labels.

Write the $g$-arrow as an arrow from the right end $y$ to the left
end $x$, and define $g(x)=y$.  If an $h$-arrow is composed on its
source side, its left end is $y$, and the composite label is $gh$.
The resulting coefficient pullback is $g\circ h$.  This convention
therefore gives left semilinear actions: the moving source coordinate
$g(x)$ supplies their coefficient action.  Strict labels are tangent
to identity, and the grading labels act as stated.  Central $[z]$
fixes $x$ and has trivial grading: its unique lift is $[z]_{F_x}$
itself.  Dividing by it gives the second presentation in
\eqref{geom:quotient-thickenings}.

For the last assertion, a strict-label lift has $\rho\in1+xR$.
There is a unique $\chi\in1+xR$ with $\chi^{1-p}=\rho$, since
$1-p$ is a unit.  The linear coordinate change $T\mapsto\chi T$
sends
\[
 v_i\longmapsto\chi^{1-p^i}v_i=\rho^{m_i}v_i.
\]
It carries the law at $y=x/\rho$ to the source of the strict
arrow.  Composing gives the required formal isomorphism over $R$.
Finally invariance of $t=xu$ gives \eqref{geom:twist-cocycle} over
the domain $R$.  Reducing this identity gives the formula on each
$R_d$; thus no division by a nilpotent coordinate is involved.
\end{proof}

\begin{proposition}\label{geom:coinduction}
The sheaves $T'$ and $T_k$ are $t$-monic.  Set
$T'_d=T'/t^d$ and $T_{k,d}=T_k/t^d$; their ordinary quotients also
represent their derived multiplication cofibers.  On the chart of
Proposition~\ref{geom:thickened-charts},
\[
 T'_d=C_{\mathrm{lc}}(G/\Lambda,R_d)
       =C_{\mathrm{lc}}(H,R_d),
\]
with action
\[
 (gF)(s)=g\bigl(F(g^{-1}s)\bigr)
       \qquad(g\in G,\ s\in G/\Lambda).
\]
The inner value is acted on through the coefficient action on $R_d$.
The analogous description of $T_{k,d}$ uses
$G_{\mathrm{up}}/\Lambda$.
\end{proposition}

\begin{proof}
The maps from $\mathcal X_k/(t^d)$ in
Proposition~\ref{geom:thickened-charts} are the subgroup maps for
$\Lambda$, so their pushforward is coinduction.  To check the
semilinear formula explicitly, let $M$ be a semilinear $G$-module.  An $R_d$-linear
$\Lambda$-equivariant functional $\epsilon:M\to R_d$ has adjoint
\[
 m\longmapsto\bigl(s\longmapsto s(\epsilon(s^{-1}m))\bigr).
\]
The function is constant on cosets, is locally constant by
smoothness, and is equivariant for the stated action.  Evaluation
at $1\Lambda$ is the inverse adjunction.  This proves the formula.
Multiplication by $t$ is injective on $\mathcal O_{\mathcal X_k}$.
Affineness and the projection formula show it remains injective
after either chart pushforward, and identify the quotients with the
pushforwards from $\mathcal X_k/(t^d)$.  Monicity makes these
ordinary quotients the multiplication cofibers of pure complexes.
\end{proof}

\subsection{The reduced stabilizer and its operation ring}

To describe the completed coefficient ring, we now study $H$ and its
action on $R$.  The formal-group digits give ordered coordinates on
$H$ and a polynomial associated graded for its completed skew group
ring.  This description also gives the cohomological dimension bound
directly, without a uniform-group hypothesis on the whole of $H$.

For $i\geq1$, let $G^1(i)$ consist of strict Honda automorphisms
whose digits $t_j$ vanish for $1\leq j<i$, and let $H(i)$ be its
image in $H$.  Set
\[
 \Gamma_s=H(2s+1)\qquad(s\geq0).
\]

\begin{lemma}\label{geom:group-filtration}
The filtration is complete and separated, and
\begin{equation}\label{geom:group-graded}
 H(i)/H(i+1)=
 \begin{cases}
 (k,+),&i\text{ odd},\\
 (k/\F_p,+),&i\text{ even},
 \end{cases}
 \qquad [H(i),H(j)]\subset H(i+j).
\end{equation}
The $p$-th-power map carries each graded digit to the identical digit
two indices later.  Choose $g_1,g_2$ with first digits an $\F_p$-basis
of $k$, and choose $g_3=\gamma$ with only nonzero digit
$t_2=\theta\notin\F_p$.  Then
\[
 (a_1,a_2,a_3)\longmapsto g_1^{a_1}g_2^{a_2}g_3^{a_3}
\]
gives unique coordinates $\Z_p^3\to H$.  Modulo $\Gamma_s$, the
coordinates are taken modulo $p^s$.  In particular $\Gamma_s$
consists of coordinates divisible by $p^s$, and $H(2)$ consists of
coordinates whose first two entries are divisible by $p$.
The element $\gamma$ commutes with $\Lambda$.
\end{lemma}

\begin{proof}
The special-fiber labels, with composition truncated on each jet,
make $G^1$ a compact Hausdorff inverse limit of finite groups.  The central subgroup $Z$ is compact and hence
closed, so $H$ is compact Hausdorff.  Each $H(i)$ is closed.  Nested
cosets in this compact group show that $H\to\varprojlim_i H/H(i)$
is surjective.  To see that its kernel is trivial, represent an
element in every $H(i)$ by $g\in G^1$.  The nonempty closed sets
\[
 \{z\in Z:g z^{-1}\in G^1(i)\}
\]
are nested.  Their intersection is nonempty, and any element in it
shows that $g\in Z$, since the digit filtration of $G^1$ is separated.
This proves completeness and separatedness of the quotient filtration.

Work first in the ring of actual endomorphisms of $F_0$, with its
formal-group addition and composition.  An automorphism beginning
in digit $i$ is $1+h$, where $h$ has order at least $p^i$.
Composition orders give the commutator inclusion.  In the binomial
expansion of $(1+h)^p-1$, the leading term at order $p^{i+2}$ is
$ph$: here $[p]_{F_0}(T)=T^{p^2}$, and for $p>3$ the remaining
terms have larger order.  Its coefficient is the original digit,
since that digit lies in $k=\F_{p^2}$.  All digits occur by the
arrow calculation above.  The central subgroup has its first
nonzero digit in $\F_p$ at an even index, as follows from
$[p^a]_{F_0}(T)=T^{p^{2a}}$.  Quotienting gives
\eqref{geom:group-graded} and the power assertion.

For the coordinates, proceed from $s$ to $s+1$.  The $p^s$-th
powers of $g_1,g_2$ correct the two independent digits of
$H(2s+1)/H(2s+2)$.  They do so also when inserted in their ordered
positions, because the resulting commutators lie in the next
filtration step.  The $p^s$-th power of $g_3$ then corrects the
remaining digit of $H(2s+2)/H(2s+3)$.  These corrections are unique.
Completeness and separatedness of the digit filtration give the
coordinates on the inverse limit and all the asserted subgroup
descriptions.  The weight $m_2=q$ is trivial on $\Lambda$ and the
other digits of $\gamma$ vanish, proving its commutation assertion.
\end{proof}

Define the completed rings, with coefficients on the left, by
\[
 \Omega=k[[H]],\qquad \Delta=R[[H]].
\]
The second is a skew group ring for the action on $R$: it is the
inverse limit of $R_d\#(H/\Gamma_s)$, with $s$ large enough to
detect the action on $R_d$.  Thus $g a=g(a)g$ for $a\in R$.
Put $U_i=g_i-1$.  The filtration below is an auxiliary filtration on
these coefficient rings; it is distinct from both $x$-adic order and
the homological grading of periodic complexes.

\begin{proposition}\label{geom:skew-filtration}
The ring $\Delta$ has a complete separated decreasing filtration
given by the ordered weights
\[
 \operatorname{wt}(x)=4,\qquad
 \operatorname{wt}(U_1)=\operatorname{wt}(U_2)=1,\qquad
 \operatorname{wt}(U_3)=3.
\]
Its associated graded is the commutative polynomial ring
\begin{equation}\label{geom:graded-operations}
 \gr\Delta=k[\sigma x,A,B,C].
\end{equation}
The element $x$ is normal, $\Omega=\Delta/(x)$, and the induced
filtration gives $\gr\Omega=k[A,B,C]$.  Here $A,B,C$ are the initial
symbols of $U_1,U_2,U_3$.
\end{proposition}

\begin{proof}
The coordinates of Lemma~\ref{geom:group-filtration} show that
ordered monomials in $U_1,U_2,U_3$ form a topological $k$-basis of
$\Omega$.  At level $H/\Gamma_s$, the ordered group elements and
the ordered difference monomials with exponents less than $p^s$
are related by an invertible binomial change of basis; larger
$p^s$-th powers vanish in characteristic $p$.  Including powers of
$x$ gives the corresponding basis of $\Delta$.

Initially filter by the closed spans of all words of weight at
least a given integer.  These spaces tend to zero.  Indeed at each
finite quotient the action preserves the ideal $(x)$, so $x$ is normal
there.  It is also nilpotent, and modulo $x$ the ideal generated by
the remaining letters is the augmentation ideal of a finite
$p$-group algebra.  Thus the ideal generated by all the letters is
nilpotent at each finite quotient.  Conversely, the
word-filtration spaces contain deep open neighborhoods by the
ordered topological basis.

There is a strict gain in weight whenever letters are interchanged.
For $U_1,U_2$, the commutator lies in $H(2)$, whose first two
coordinates are divisible by $p$.  Characteristic-$p$ binomial
expansion therefore gives weight at least $3$ for its difference
from 1, exceeding the original weight $2$.  Commuting $U_1$ or
$U_2$ with $U_3$ gives an element of $\Gamma_1$, so the difference
has weight at least $p>4$.  Finally
\[
 U_i x-xU_i=(g_i(x)-x)g_i
\]
has weight at least $8$, since $g_i(x)-x\in x^2R$.  This is larger
than either $\operatorname{wt}(U_i)+4$.

We can consequently reorder every word modulo strictly higher
weight, and iterate to all weights by completeness.  Every element
of a word-filtration stage expands in ordered monomials of at least
that weight.  Uniqueness in the topological basis proves the reverse
inclusion and shows that there are no further associated-graded
relations.  This proves \eqref{geom:graded-operations}.

Normality of $x$ follows from
$g x=x(g(x)/x)g$.  Weighting each group coefficient by the continuous
unit $g(x)/x$, or by its inverse, defines the requisite continuous
change of coefficients even in the completion: at precision $x^d$
choose a group quotient detecting the action modulo $x^{d+1}$.
Thus the ideal $(x)$ is two-sided and quotienting simply removes
the $x$ coordinate and its weight.
\end{proof}

\subsection{Derived sections and cohomological bounds}

We now turn the chart calculations into bounds for derived sections.
First we justify the model and resolution tools for the periodic
categories.  The polynomial associated graded of $k[[H]]$ will yield
a compact resolution of $k$ of length three.  Applying continuous Hom
computes derived $H$-invariants; exact $\Lambda$-invariants and the
forgetful functor from semilinear $R_d$-modules then give the bound on
the nilpotent charts.  Passing through the open, the central quotient,
and the integral base adds the degrees recorded below.  The bounds are
uniform for arbitrary sheaves, so they also control unbounded periodic
complexes.

\begin{lemma}\label{geom:injective-model}
For each flat affine groupoid stack used here, twisted periodic
complexes admit the injective model structure with quasi-isomorphisms
as weak equivalences and monomorphisms as cofibrations.  It is stable,
its fibrant objects have injective sheaf terms, and maps into a
fibrant object are computed up to compatible chain homotopy.
If a sheaf $W$ satisfies $\Ext^{>a}(W,-)=0$ for some finite $a$, maps
from its pure shifts can be computed by applying $\Hom(W,-)$
termwise to any quasi-isomorphic replacement whose terms are
$\Hom(W,-)$-acyclic.
\end{lemma}

\begin{proof}
The sheaf categories in question are Grothendieck categories.  Colimits
and AB5 follow from flat affine descent.  For a generating set, choose
an infinite cardinal $\kappa$ dominating the sizes of the base and
arrow rings.  Starting with a bounded-size family of elements of a
comodule, repeatedly adjoin components in finite expressions for
their coactions and close under the module operations.  Flatness
injects the tensor of the resulting submodule into that of the
original module, so the coaction identities restrict.  The resulting
subcomodules have bounded cardinality and form a generating set.
For $\mathcal X$ one may use its finitely many grading-trivialized
principal charts; their disjoint union is affine and the same argument
applies.  The other stacks have the affine presentations constructed
above.

Twisted periodic complexes are Grothendieck as well.  Disk objects
formed from sheaf generators in adjacent indices, with identity
differential and then repeated with the twists, give generators.
We use Smith's recognition theorem in the form of
\cite[Theorem~1.7]{Beke2000}; the solution-set condition follows from
accessibility by \cite[Proposition~1.15]{Beke2000}.
Quasi-isomorphisms satisfy the two-out-of-three property and are closed
under retracts, as is seen on homology.  They are accessibly embedded
and accessible.  For an explicit size argument,
enlarge $\kappa$ to include the twist data.  In an arrow of complexes
start with boundedly many elements, close under coactions and
differentials, and then adjoin preimages for cycles in its cone.
Iteration expresses any quasi-isomorphism as a $\kappa^+$-directed
union of quasi-isomorphic subarrows of size at most $\kappa$.
Objects and arrows of that size are $\kappa^+$-presentable: their
module presentations, coaction identities, and finitely many
periodic differential identities can all be tested at that size.
These tests can be performed on the flat affine charts.

Acyclic monomorphisms are preserved by pushout and transfinite
composition, by AB5.  Monomorphisms are cofibrantly generated: for a
generator $G_0$, adjoin successively the images of maps from $G_0$ by
pushing out inclusions of its subobjects, taking the inverse image
of the subobject already attained at each step.  There is a set of
these subobject inclusions.  A map right orthogonal to all
monomorphisms is a split epimorphism with injective kernel.  An
injective object in the category of complexes is contractible:
its monomorphism into the cone of its identity splits, making it a
retract of a contractible complex.  Such an epimorphism is therefore
a quasi-isomorphism.  This verifies the remaining recognition
condition.

Ordinary chain mapping cylinders are compatible with the twists;
they identify ordinary mapping cones with homotopy cofibers.
The chain shift is an equivalence, so the model structure is stable.
Testing fibrancy against disk inclusions shows that fibrant complexes
have injective terms.  Since all objects are cofibrant, the chain
cylinder also gives the claimed computation of maps by chain homotopy.

For the last assertion, a fibrant replacement computes the derived
maps.  A quasi-isomorphism between complexes with
$\Hom(W,-)$-acyclic terms remains a quasi-isomorphism after applying
$\Hom(W,-)$.  Indeed its cone is acyclic with such terms.  Dimension
shifting along the short exact sequences of cycles shifts the
possible obstruction to exactness past degree $a$, where it vanishes.
The argument requires only the uniform bound on $\Ext(W,-)$ and
therefore works for the unbounded periodic unfolding.
\end{proof}

\begin{lemma}\label{geom:sections-resolution}
The right derived functors of sections in the sheaf categories above
commute with filtered colimits and are computed by the flat affine coinduced
resolution.  The resolution is compatible with closed base change.
If $\Ext^{>a}(\mathcal O,-)=0$, its finite truncation gives a
sections-acyclic resolution that computes \eqref{geom:sections-homotopy}
on periodic complexes.  In particular $\pi_*$ then commutes with
filtered colimits and telescopes.
\end{lemma}

\begin{proof}
Choose an affine fpqc chart $l:S\to\mathcal S$ with $S$ affine.
The functor $I=l_*l^*$ is exact, its unit is monic, and its unit
cokernel $J$ is exact.  Its values are sections-acyclic by flat
affine adjunction.  Iterating $I$ and $J$ gives the usual coinduced
resolution.  On the chart all these functors, including ordinary
sections of the coinduced terms, commute with filtered colimits;
AB5 therefore proves the same for their cohomology.  After restricting
to a closed substack, the base-changed chart gives exactly the
corresponding resolution.  This proves the assertion about sheaves
pushed forward from closed loci without requiring derived closed
pushforward to be fully faithful.

Under the dimension bound, the resolution can end in $J^a$, which is
sections-acyclic by dimension shifting.  These functors commute with
weight twists.  The finite totalization thus computes sections on
periodic complexes, by Lemma~\ref{geom:injective-model}; no infinite
totalization is needed.  The colimit assertion follows termwise.
\end{proof}

We use exact pullback and pushforward functors below also on periodic
complexes.  An exact pullback left adjoint is left Quillen for these
injective structures and is adjoint to the right derived pushforward.
These adjunctions let the preceding chart calculations control both
sheaves and the periodic complexes built from them.

\begin{lemma}
Invariants\label{geom:group-cohomology} of smooth $k$-linear $H$-modules have cohomological
dimension at most $3$.  The same is true for $G$ and for sections
on $\mathcal N/(t^d)$, for every $d\geq1$.
\end{lemma}

\begin{proof}
The key step is to lift the three-step Koszul resolution from the
associated graded ring and use it to compute continuous invariants.
Lift the graded Koszul resolution of $k$ over
$\gr\Omega=k[A,B,C]$ to a resolution of length three by finite free
compact left $\Omega$-modules.  Here is the lifting argument.  At
each step the actual kernel is complete.  Lift the homogeneous
columns generating its associated-graded kernel to elements of that
kernel and map in a free module with the corresponding shifted
filtration.  Successive subtraction of initial symbols converges
and makes this a strict surjection.  Its complete kernel has
associated graded equal to the next Koszul kernel.  At the final
step that graded kernel vanishes, so separatedness makes the actual
kernel zero.  The resulting differentials compose to zero because
their columns were lifted into the actual kernels.

For a smooth discrete module $M$, continuous
$\Hom_\Omega(-,M)$ applied to this resolution has terms finite sums
of $M$, and hence gives a delta functor.  It is acyclic in positive
degrees on $C_{\mathrm{lc}}(H,M_{\mathrm{vec}})$, coinduced from the
underlying vector space: evaluation and translation identify the
resulting complex with continuous $k$-linear Hom from the compact
resolution into $M_{\mathrm{vec}}$.  The adjunction is continuous
because a map from a compact free module to a discrete vector space
kills a sufficiently deep stable filtration stage.  Exact sequences
of compact $k$-vector spaces split continuously, as is seen by
dualizing to discrete vector spaces; thus this Hom complex is exact
in positive degrees.  Every smooth module embeds in such a
coinduction.  Effaceability identifies the delta functor with derived
invariants and proves the bound.

For $G$, take $H$-invariants and then exact $\Lambda$-invariants.
Restriction of injectives from $G$ to $H$ preserves injectives,
because finite-index induction is exact.  Finally, forgetting the
$R_d$-module structure of a semilinear representation preserves
injectives: it has the exact left adjoint given by the free
$R_d$-module with diagonal action.  Applying the quotient-chart
description proves the assertion about $\mathcal N/(t^d)$.
\end{proof}

\begin{proposition}\label{geom:cohomological-bounds}
Projective dimensions in the Ext sense satisfy
\begin{equation}\label{geom:dimension-bounds}
 \operatorname{pd}_{\mathcal N}(\mathcal O)\leq4,\qquad
 \operatorname{pd}_{\mathcal M_k}(\mathcal O),
 \operatorname{pd}_{\mathcal M_0}(\mathcal O)\leq5,\qquad
 \operatorname{pd}_{\mathcal M}(\mathcal O)\leq6.
\end{equation}
For sheaves on $\mathcal M_k$, the right adjoint of $r^*$ has
\begin{equation}\label{geom:central-derived-invariants}
 r_*M=\ker(c-1:M\to M),\qquad
 R^1r_*M=\coker(c-1:M\to M),\qquad R^{>1}r_*M=0,
\end{equation}
where the kernel and cokernel descend to $\mathcal N$.
\end{proposition}

\begin{proof}
We first pass from the nilpotent charts to $\mathcal N$, then restore
the central group, and finally pass from characteristic $p$ to the
integral stack.

A $t^\infty$-torsion sheaf on $\mathcal N$ is a filtered union of
sheaves annihilated by powers of $t$, with the evident weight twists
in the kernel maps.  Lemmas~\ref{geom:sections-resolution}
and~\ref{geom:group-cohomology} show that its sections cohomology
vanishes above degree $3$.  The open of
Proposition~\ref{geom:open} has exact sections.  For an arbitrary
sheaf $M$, the kernel and cokernel of its localization map
$M\to M[1/t]$ are $t^\infty$-torsion.  The open is affine over the
stack, so localization is exact open restriction followed by affine
pushforward.  The two short exact sequences associated to the kernel
and cokernel give the extra degree and prove the bound $4$.

The right adjoint $r_*$ is vertical invariants.  Smoothness and
centrality ensure that the kernel and cokernel of $c-1$ have trivial
vertical action and descend; they give a two-term delta functor.
To prove that it computes the derived functor, efface its degree-one
term by embedding a sheaf into $I=l_*l^*$ for an affine law chart
$l:S\to\mathcal M_k$.  On the arrow chart, the action is translation
on the arrow algebra tensored over the other endpoint with the
module coefficients.  The map $c-1$ is surjective there.  Indeed the
arrow scheme is a $Z$-torsor over its quotient, and fpqc locally it
is split.  A locally constant function is represented at a finite
central period.  Increase that period by a factor $p$; the sum of
the repeated old period is zero in characteristic $p$, so the
function can be integrated as a discrete difference on the larger
cycle.  This proves surjectivity locally, hence on the sheaf.
Effaceability proves \eqref{geom:central-derived-invariants}.

Because $r^*$ is exact, $r_*$ preserves injectives.  Derived sections
composed with \eqref{geom:central-derived-invariants} therefore give
the bound $5$ on $\mathcal M_k$.  Finite flat extension from $\F_p$
to $k$ commutes with the coinduced chart resolution of sections and
detects vanishing, giving the same bound on $\mathcal M_0$.

For $\mathcal M$, first consider $p$-power-torsion sheaves.  A sheaf
annihilated by $p^a$ has a finite filtration with characteristic-$p$
subquotients; a general $p^\infty$-torsion sheaf is the filtered union
of its kernels of powers of $p$.  The preceding bound and
Lemma~\ref{geom:sections-resolution} give vanishing above degree
$5$.  On the rational open the logarithm trivializes the strict law
groupoid, leaving only $B\Gm$, whose sections are exact.  Apply the
same kernel--cokernel argument to $M\to M[1/p]$ to obtain the last
bound in \eqref{geom:dimension-bounds}.
\end{proof}

\subsection{The central comparison triangle}

The finite resolutions above justify derived sections on periodic
complexes.  For the reverse transfer, we now construct the central
pushforward itself from exact functors and retain the natural comparison
triangle, including its possibly nonzero connecting map.

\begin{proposition}\label{geom:relative-pushforward}
Choose an affine law chart $l:S\to\mathcal M_k$, set
$I=l_*l^*$, and let $J$ be the cokernel of its unit.  The two-term
totalization
\begin{equation}\label{geom:Q-definition}
 Q=\bigl[r_*I\longrightarrow r_*J\bigr],
 \qquad\text{in homological degrees }0,-1,
\end{equation}
computes $Rr_*$ on twisted periodic complexes and preserves
quasi-isomorphisms.  There are natural identifications
\[
 QT_k\simeq T',\qquad
 \pi_*(r^*U)\cong\pi_*(Qr^*U)\quad(U\in\mathbf D(\mathcal N)),
\]
and natural triangles
\begin{equation}\label{geom:central-triangle}
 U\longrightarrow Qr^*U\longrightarrow U[-1]
                     \longrightarrow U[1].
\end{equation}
\end{proposition}

\begin{proof}
The composite $rl:S\to\mathcal N$ is an affine chart, so $I$ has
$r_*$-acyclic values by affine adjunction.  Equation
\eqref{geom:central-derived-invariants} and the sequence
$0\to M\to IM\to JM\to0$ show that $J$ also has $r_*$-acyclic
values.  Both $r_*I$ and $r_*J$ are exact: $I$ and $J$ are exact,
and the acyclicity of their values removes the connecting
obstruction after applying $r_*$.  Hence the finite totalization
\eqref{geom:Q-definition} preserves quasi-isomorphisms.  On
injective-term fibrant replacements it agrees up to quasi-isomorphism
with ordinary $r_*$, by the same two-term acyclic resolution.
It therefore computes $Rr_*$ on periodic complexes, not just on
sheaves.  The twist line is pulled back from downstairs, so all
these functors and identifications respect the periodicity.

Flat affine adjunction along $f_k$ and $f'$ identifies
$Rr_*T_k$ with $T'$.  Derived adjunction for $r^*$, with
$r^*\mathcal O=\mathcal O$, gives the equality of sections homotopy.
For a sheaf pulled back from $\mathcal N$, the vertical kernel and
cokernel in \eqref{geom:central-derived-invariants} are both the
original sheaf.  Thus there is a natural four-term exact sequence
\[
 0\longrightarrow U\longrightarrow r_*Ir^*U
 \longrightarrow r_*Jr^*U\longrightarrow U\longrightarrow0.
\]
It applies termwise also to periodic complexes.  Its two connecting
maps give \eqref{geom:central-triangle}, with the consistent
homological signs.  In particular the triangle is natural, but
no splitting of its final connecting map is asserted.
\end{proof}

The next sections carry out the algebraic construction in
$\mathbf D(\mathcal N)$.  The open is $\Spec k$, the thickenings carry
the semilinear $G$-action, and the reductions of $T'$ are coinduced
modules.  Section~\ref{complete:section} uses these descriptions and
the finite bounds to assemble the reductions by homotopy limits.
